\chapter{Conclusion and Outlook}
\label{chap:conclusion}

This thesis examined 12-month OSS dependency inactivity prediction from public
repository health and maintenance indicators, together with parallel
security-practice context for triage. On the time-aware held-out split,
full-history XGBoost reaches ROC-AUC \(0.958\), Brier \(0.078\), and expected
calibration error \(0.035\). ROC-AUC remains \(0.950\) under a
repository-disjoint split, bounding optimism from repository recurrence.

The aggregate result mostly reflects recency. The last-commit rule reaches
\(0.957\) on its \(4{,}556\) defined rows; XGBoost reaches \(0.974\) on the same
rows (Section~\ref{sec:results-baselines}). The full test set has inactive rate
\(0.661\) and contains many repositories already quiet at \(t\). On
active-at-\(t\), with inactive rate \(0.225\), XGBoost reaches \(0.905\) against
\(0.844\) for the best trivial rule. The margin is \(0.061\), compared with
\(0.016\) on recency-defined rows.

Public repository signals therefore rank 12-month inactivity accurately within
the examined dataset, but recency provides most aggregate predictive content.
Combining signals adds a modest, measurable improvement for repositories still
active at the observation date.

Performance falls to \(0.846\) on low-visibility repositories, and recall on
active-at-\(t\) is \(0.49\) at the selected threshold. Probabilities are
calibrated to sampled prevalence. The score therefore prioritizes review and
reports per-prediction evidence support; it does not determine abandonment.
Security-practice signals remain separate context and are excluded from the
inactivity model (Section~\ref{sec:related-security-practice}). The practical
contribution is the per-repository ranked view in
Chapter~\ref{chap:demonstrator}.

\section{Outlook and Future Work}
\label{sec:conclusion-outlook}

The remaining work is:

\begin{itemize}
    \item \textbf{External validation.} The time-aware split remains inside the foundation dataset (Section~\ref{sec:method-evaluation-strategy}). Evaluation on a fresh repository sample would test generalization beyond its prevalence and seed buckets.
    \item \textbf{A frame sampled at the observation date.} The foundation seed was drawn from repository state at generation time, so entry into the frame depends on behavior inside the outcome windows and all reported metrics are conditional on that construction (Sections~\ref{sec:method-retrospective-frame} and~\ref{sec:validity-external}). Building the frame from a repository listing that predates each observation date, so that inclusion uses only information available at \(t\), would remove the dependency and give an aggregate figure comparable to deployment.
    \item \textbf{Systematic qualitative error analysis.} Active-at-\(t\) and the trivial baselines isolate predictive skill from persistence (Sections~\ref{sec:results-slices} and~\ref{sec:results-baselines}). Larger-scale case review is still needed to find predictions made for the wrong reason.
    \item \textbf{Evaluation with analysts.} The demonstrator separates a small candidate set and shows divergent inactivity and security-posture signals, but it does not establish better analyst decisions or throughput (Chapter~\ref{chap:demonstrator}). A controlled practitioner study should compare review time, agreement, and decision quality with and without the ranking and parallel posture column. Public repository data alone cannot answer this question.
    \item \textbf{Additional model families.} More expressive deep-learning architectures and sequence models over repository event history remain untested.
    \item \textbf{Ensemble evaluation.} Runtime averages the staged models of a regime, while reported metrics cover individual artifacts (Section~\ref{sec:results-performance}). Evaluation should compare the ensemble with its strongest member and test model-agreement spread as an uncertainty signal.
    \item \textbf{Prevalence-aware calibration.} The seed oversamples dormant and archived repositories, so probabilities reflect sampled prevalence (Section~\ref{sec:method-limitations}). Recalibration to deployment prevalence is required for live probabilities.
    \item \textbf{Extended label sensitivity.} Section~\ref{sec:results-label-sensitivity} varies two-of-four to one-of-four and three-of-four. Individual active-month, contributor, merged-pull-request, and release thresholds remain to be swept. The mature libraries in Section~\ref{sec:demo-ranking} show the underlying construct problem: public signals may not distinguish stable maturity from decline (Section~\ref{sec:validity-construct}).\kiref{Claude Opus 5/8}
    \item \textbf{Cold-start coverage and on-demand history reconstruction.} Cold-start applies outside the seed, uses one live snapshot, and omits historical signals, including annual commit distribution. Its metrics are measured on seed repositories but applied to arbitrary submissions, so practical performance may be optimistic; evidence support marks such cases. Recent history could be reconstructed through the GitHub REST API or GH Archive's public BigQuery dataset. The former offers weekly commit and contributor statistics; the latter avoids per-repository rate limits but adds an external query dependency. Neither supplies the full feature set within one request. First-response latency requires prohibitively many issue timelines, and year-over-year features require a second year. Imputing this missing tail into full-history would introduce distribution shift. A third \emph{warm-start} model should instead use only reliably reconstructed live-API features, be validated against the seed cache, and use a rate-limit-aware request cache. This requires further engineering and validation.
\end{itemize}
